\section{贡献与评价}
\subsection{5.1 主要发现}
\begin{frame}{数据为模型的信息输入提供了哪些约束？}
  \strap{5.1\quad 回到开场：投资者实际获得了什么}
  {\small\begin{tabularx}{\linewidth}{@{}L{.13\linewidth}YL{.19\linewidth}@{}}
    \toprule
    \textbf{对象} & \textbf{从数据看到的事实} & \textbf{证据}\\
    \midrule
    投入 & 典型家庭对已交易股票的可见研究有限，分布右偏。 & 表2\\
    内容 & 价格图覆盖广。默认展示和主动调整对应不同历史窗口。 & 表3、图1\\
    差异 & 股票覆盖与内容侧重有别。家庭差异体现于参与程度和主题。 & 表4–5\\
    时点 & 同股研究在交易附近集中。许多路径在一个主题后结束。 & 图3–5\\
    \bottomrule
  \end{tabularx}}
  \finding{理论可对照具体的内容、期限和获取时点，再研究它们怎样进入预期与需求。}
  \refnote{作者贡献与研究方向见论文 pp.~5–7、40–41。}
\end{frame}

\subsection{5.2 评价与后续}
\begin{frame}{页面访问可直接看见，内容占比依赖分配}
  \strap{5.2\quad 汇报人评价}
  \begin{columns}[T,onlytextwidth]
    \column{.46\textwidth}
      \shortlabel{直接分类}\par
      \greatnum{93\%}\par
      {\small 家庭访问过价格图。约1\%专门查看传统风险统计。}
    \column{.49\textwidth}
      \shortlabel{内容归因}\par
      \greatnum{\(\approx 42\%\)}\par
      {\small 均值研究分钟来自三类混合页。分配后，价格图约占44.7\%。}
  \end{columns}
  \medskip
  \finding{进一步比较不同分配权重下，主题排序是否稳定，份额变化有多大。}
  \refnote{Table 3（p.~48）、附录 pp.~66–68、Table IA2。42\%和44.7\%据表计算。}
\end{frame}

\begin{frame}{哪些实际接触的信息改变了判断？}
  \strap{5.2\quad 以默认窗口为例的后续检验}
  \begin{center}
    \begin{tikzpicture}
      \node[storyfocus,text width=2.7cm] (a) {随机改变\\默认回看窗口};
      \node[storybox,text width=2.7cm,right=6mm of a] (b) {记录实际曝光\\与主动调整};
      \node[storybox,text width=2.7cm,right=6mm of b] (c) {比较前后预期\\与后续选择};
      \draw[grapharrow] (a)--(b);
      \draw[storyarrow] (b)--(c);
    \end{tikzpicture}
  \end{center}
  \vspace{3mm}
  \finding{本文把可得资料与个体实际获取的内容区分开来。现在可以进一步检验其中哪些信号进入判断和行动。}
  \refnote{默认窗口实验为汇报人拟议设计。论文的后续问题见 p.~41。}
\end{frame}
